% ============================================================
% Section 122: Ge、Si 半导体价带简并与轻重空穴
% ============================================================
\section{半导体物理：Ge、Si 价带简并与轻重空穴}
\label{sec:ge-si-valence-band}

\begin{modelbox}[题目：Ge、Si 半导体价带简并分析]
Ge、Si 半导体价带简并，价带顶附近电子能量表达式为
\[
E_{1,2}(k)=E_V-\frac{\hbar^2}{2m}\left\{Ak^2\pm\left[B^2k^4+C^2(k_x^2k_y^2+k_y^2k_z^2+k_z^2k_x^2)\right]^{1/2}\right\},
\]
其中 $m$ 为电子质量，$A,B,C$ 为常数。

(1) 用球形等能面近似，求出轻、重空穴的有效质量，在电子能带图上，示意画出并标明轻、重空穴带；

(2) 比较相同波矢轻、重空穴能量，并在空穴能带上标出轻、重空穴带；

(3) 在一定温度下，重空穴数目多还是轻空穴多，为什么？
\end{modelbox}

\begin{tipbox}[考点快速分析]
\begin{itemize}
    \item \textbf{球形等能面近似}：忽略各向异性项，令 $C=0$
    \item \textbf{轻重空穴来源}：价带顶在 $\Gamma$ 点简并，自旋-轨道耦合分裂为两支
    \item \textbf{空穴有效质量}：$m_h^*=-m_e^*$，曲率越大（带越陡）空穴质量越小
    \item \textbf{载流子浓度}：态密度 $g(E)\propto (m^*)^{3/2}$，有效质量大则态密度大
\end{itemize}
\end{tipbox}

\begin{derivationbox}[标准解题过程]
\textbf{(1) 球形等能面近似下的轻重空穴有效质量}

在球形等能面近似下，令 $C=0$，价带顶附近的两支能带简化为：
\[
E_1(k)=E_V-\frac{\hbar^2}{2m}(A+B)k^2,\qquad
E_2(k)=E_V-\frac{\hbar^2}{2m}(A-B)k^2,
\]
其中 $A>B>0$。两支能带均呈抛物线形，但曲率不同。

\textit{电子有效质量}：
\[
m^*=\hbar^2\left(\dd[2]{E}{k}\right)^{-1}.
\]

$E_1$ 支：$\dd[2]{E_1}{k}=-\dfrac{\hbar^2}{m}(A+B)$，故 $m_{e1}^*=-\dfrac{m}{A+B}$。

$E_2$ 支：$\dd[2]{E_2}{k}=-\dfrac{\hbar^2}{m}(A-B)$，故 $m_{e2}^*=-\dfrac{m}{A-B}$。

\textit{空穴有效质量}：$m_h^*=-m_e^*$。

轻空穴（对应 $E_1$ 支，曲率大、质量小）：
\begin{equation}
\boxed{m_{lh}^*=\frac{m}{A+B}.}
\end{equation}

重空穴（对应 $E_2$ 支，曲率小、质量大）：
\begin{equation}
\boxed{m_{hh}^*=\frac{m}{A-B}.}
\end{equation}

在电子能带图（纵轴为电子能量，上方为导带，下方为价带）中：价带顶在 $E_V$ 处。$E_1$ 支曲率较大（较陡），为轻空穴带；$E_2$ 支曲率较小（较平），为重空穴带。

\textbf{(2) 相同波矢下轻、重空穴的能量比较}

空穴的波矢 $k_h=-k_e$，空穴能量（以价带顶为零点，向下为正）为：
\[
E_{h1}(k_h)=\frac{\hbar^2}{2m}(A+B)k_h^2,\qquad
E_{h2}(k_h)=\frac{\hbar^2}{2m}(A-B)k_h^2.
\]
由于 $A+B>A-B>0$，对于相同的波矢大小 $k_h$：
\begin{equation}
\boxed{E_{h1}(k_h)>E_{h2}(k_h),}
\end{equation}
即轻空穴的能量比重空穴更高。在空穴能带图（纵轴为空穴能量，向下为正）中，轻空穴带在上方，重空穴带在下方。

\textbf{(3) 一定温度下轻重空穴的数密度比较}

空穴服从费米-狄拉克分布。在热平衡下，能量越低的状态被占据的概率越大。由于重空穴带在相同波矢下能量更低，因此重空穴态更容易被空穴占据。

定量地，三维情况下价带顶附近态密度有效质量 $m_d^*\propto (m^*)^{3/2}$。由于 $m_{hh}^*>m_{lh}^*$，重空穴的态密度远大于轻空穴。因此，在任意有限温度下：
\begin{equation}
\boxed{\text{重空穴的数目远多于轻空穴。}}
\end{equation}
通常所说的``空穴''在多数输运过程中指重空穴，因其占绝对多数。
\end{derivationbox}

\begin{formulabox}[答案汇总]
\begin{center}
\begin{tabular}{ll}
\toprule
\textbf{问题} & \textbf{答案} \\
\midrule
(1) 轻空穴有效质量 $m_{lh}^*$ & $\dfrac{m}{A+B}$（较小） \\
(1) 重空穴有效质量 $m_{hh}^*$ & $\dfrac{m}{A-B}$（较大） \\
(2) 相同 $k_h$ 的能量 & 轻空穴能量 $>$ 重空穴能量 \\
(3) 数目比较 & 重空穴数目远多于轻空穴 \\
\bottomrule
\end{tabular}
\end{center}
\end{formulabox}

\begin{insightbox}[物理图像]
\begin{itemize}
    \item Ge、Si 的价带顶在 $\Gamma$ 点发生简并，自旋-轨道耦合使其分裂为两支：轻空穴带和重空穴带。两者的有效质量不同，源于能带曲率的差异。
    \item 重空穴有效质量大，能带平缓，态密度大；轻空穴有效质量小，能带陡峭，态密度小。
    \item 热平衡下，空穴优先占据能量较低的重空穴带，因此重空穴是多数载流子，主导了空穴的输运性质。在分析 p 型半导体时，通常主要考虑重空穴的贡献。
\end{itemize}
\end{insightbox}
